%documentclass[a4paper,landscape]{article}.txt
%\documentclass[a4paper,landscape]{article}
%\usepackage[margin=1cm]{geometry}
%\usepackage{tikz}

\usetikzlibrary{shapes.geometric, arrows, positioning, calc}

%\begin{document}

\pagestyle{empty}

\tikzset{
    input/.style = {rectangle, draw, fill=orange!30, text centered, minimum height=0.7cm, minimum width=2.2cm, rounded corners=2pt, align=center, font=\small\sffamily},
    tool/.style  = {rectangle, draw, fill=blue!30, text centered, minimum height=0.7cm, minimum width=2.2cm, rounded corners=2pt, align=center, font=\small\sffamily},
    dataset/.style = {rectangle, draw, fill=green!30, text centered, minimum height=0.7cm, minimum width=2.2cm, rounded corners=0pt, align=center, font=\small\sffamily},
    highlight/.style = {dataset, draw=red!90!black, thick, font=\small\sffamily\bfseries},  % <-- nouveau style
    arrow/.style = {thick, -stealth},
}

\begin{tikzpicture}[node distance=0.8cm and 1.5cm]

% ---- Entrées ----
\node[input] (A) {S-1 ESA products};
\node[input, below=0.8cm of A] (B) {APIs: CDSE, \\ EODMSRapids, \\ PODAAC SWOT};
\node[input, below=0.8cm of B] (C) {reference products\\ at Ifremer};

% ---- Outil de construction des matchups ----
\node[tool, right=1.5cm of A] (build_match) {tool to \\ build matchups};

\draw[arrow] (A) -- (build_match);
\draw[arrow] (B) -- (build_match);
\draw[arrow] (C) -- (build_match);
% ---- Dataset matchups ----
\node[dataset, right=1.8cm of build_match] (matchups) {matchups files};

\draw[arrow] (build_match) -- (matchups);

% ---- Branche 1 : co-aligned -> split -> test (à gauche) ----
% Départ diagonal vers le bas-gauche (angle 225°)
\node[tool, below left=2.2cm and 1.0cm of matchups] (coalign) {tool to \\ build co-aligned};
\node[highlight, below=0.8cm of coalign] (coaligned_dataset) {co-aligned parquet};  % <-- style highlight
\node[tool, below=0.8cm of coaligned_dataset] (split) {tool to \\ split co-aligned};
\node[dataset, below=0.8cm of split] (splitout) {splited listings};
\node[tool, below=0.8cm of splitout] (gentest) {tool to \\ build test};
\node[dataset, below=0.8cm of gentest] (test) {test dataset};

\draw[arrow] (matchups.south) -- (coalign.north); % diagonale directe

\draw[arrow] (coalign) -- (coaligned_dataset);
\draw[arrow] (coaligned_dataset) -- (split);
\draw[arrow] (split) -- (splitout);
\draw[arrow] (splitout) -- (gentest);
\draw[arrow] (gentest) -- (test);

% ---- Branche 2 : co-locations -> training -> train & val (à droite) ----
% Départ diagonal vers le bas-droit (angle 315°)
\node[tool, below right=2.2cm and 1.0cm of matchups] (coloc_builder) {tool to \\ build co-locations};
\node[dataset, below=0.8cm of coloc_builder] (coloc) {co-locations};
\node[tool, below=0.8cm of coloc] (train_builder) {tool to \\ build training};
\node[dataset, below left=0.6cm and 0.5cm of train_builder] (train) {train dataset};
\node[dataset, below right=0.6cm and 0.5cm of train_builder] (val) {val dataset};

\draw[arrow] (matchups.south) -- (coloc_builder.north);

\draw[arrow] (coloc_builder) -- (coloc);
\draw[arrow] (coloc) -- (train_builder);
\draw[arrow] (train_builder) -| (train);
\draw[arrow] (train_builder) -| (val);

% ---- Jonction entre les branches : lien de splitout vers coloc_builder ----
\draw[arrow] (splitout) -- (coloc_builder);

\end{tikzpicture}

%\end{document}